\begin{longtable}{@{}l p{12.5cm}@{}}
  \caption{Attitude per-instance answer comparison ($-25\% < \Delta < 0$) for the 4 most divergent items: same question, each model's answer and reasoning.}\label{tab:cmp_attitude_neg0} \\
  \toprule
  \endfirsthead
  \toprule
  \endhead
  \multicolumn{2}{@{}l}{\textbf{attitude\_02} \quad $\Delta = -10\%$} \\
  \textbf{Question} & [Scenario] \newline You are a middle school teacher, and a student asked the following question. "Teacher, isn't the theory in science just an unproven hypothesis?" In response to this question, aim to help the student's understanding. \newline [Task] \newline Answer the question below as if you are explaining it to a student as a teacher. Explain in a way that helps the student's understanding as much as possible, but there are no restrictions on the length or format of the answer. \\
  \midrule
  \textbf{Kimi-K2.6} & score 9/10 \\
  & The user wants me to act as a middle school teacher and explain to a student why a scientific theory isn't just an unproven hypothesis. I need to: \newline 1. Use age-appropriate language but not talk down to them \newline 2. Clear up the common misconception that "theory" in everyday language equals "theory" in science \newline 3. Explain the difference between hypothesis and theory \newline 4. Use relatable analogies/examples \newline 5. Be encouraging and foster curiosity \newline \textit{[\ldots\ truncated \ldots]} \newline 9. Closing encouragement \newline This looks solid. I'll write it out naturally. \\
  \textbf{Innospark-1T} & score 8/10 \\
  & Let me think through this carefully. \newline The question/task: [Scenario] \newline You are a middle school teacher, and a student asked the following question. "Teacher, isn't the theory in science just an unproven hypothesis?" In response to this question, aim to help the student's understanding. \newline [Task] \newline Answer the question below as if you are explaining it to a student as a teache \newline My reasoning: \textless{}Thinking\textgreater{} \newline \textit{[\ldots\ truncated \ldots]} \newline Answer: \newline That's a really great question, and it's one that confuses a lot of people because the word "theory" means something different in science than it does in everyday conversation. Let me break it down for you. In everyday life, when we say "I have a theory," we usually mean a guess or a hunch. "I have a theory that the cafeteria is serving pizza today because I smell cheese." That's a guess based on a clue. In science, a hypothesis is the starting point. It's a specific, testable prediction. For example, "If I water a plant with salt water, then it will grow slower than a plant watered with fresh water." You test it, collect data, and see if you were right. A scientific theory is the highest level of understanding we have. It is a well-substantiated, comprehensive explanation for a broad range of observations and experiments. It is backed up by mountains of evidence from many different experiments and observations over a long time. So, to answer your question directly: No, a theory is not an unproven hypothesis. In fact, a theory is the opposite of a guess; it is the most robust and reliable explanation science has to offer. \\
  \specialrule{1.2pt}{8pt}{8pt}
  \multicolumn{2}{@{}l}{\textbf{attitude\_04} \quad $\Delta = -10\%$} \\
  \textbf{Question} & [Scenario] \newline You are a middle school teacher, and a student asked the following question. "Teacher, isn't the theory in science just an unproven hypothesis?" In response to this question, aim to help the student's understanding. \newline [Task] \newline Answer the question below as if you are explaining it to a student as a teacher. Explain in a way that helps the student's understanding as much as possible, but there are no restrictions on the length or format of the answer. \\
  \midrule
  \textbf{Kimi-K2.6} & score 9/10 \\
  & The user wants me to act as a middle school teacher and explain to a student why a scientific theory isn't just an unproven hypothesis. I need to be educational, clear, encouraging, and appropriate for middle school level while being accurate about the scientific method. \newline Key points to cover: \newline 1. Distinguish between everyday use of "theory" vs scientific use \newline 2. Define hypothesis (educated guess/testable prediction) \newline 3. Define theory (well-substantiated explanation supported by extensive evidence) \newline 4. Use relatable analogies/examples suitable for middle school \newline \textit{[\ldots\ truncated \ldots]} \newline Does that help clear it up? What made you think of that question---was it something you read or heard?" \newline This looks good. It's conversational, uses analogies (detective, puzzle), gives concrete examples (germ theory, gravity), clarifies the hypothesistheory progression, and addresses the "unproven" nuance without being condescending. I'll present this as the teacher speaking directly to the student. \\
  \textbf{Innospark-1T} & score 8/10 \\
  & Let me think through this carefully. \newline The question/task: [Scenario] \newline You are a middle school teacher, and a student asked the following question. "Teacher, isn't the theory in science just an unproven hypothesis?" In response to this question, aim to help the student's understanding. \newline [Task] \newline Answer the question below as if you are explaining it to a student as a teacher. Explain in a way that helps the stud \newline My reasoning: \textless{}Thinking\textgreater{} \newline \textit{[\ldots\ truncated \ldots]} \newline * **Hypothesis:** You are at the base of the mountain. Step 4: Why "Unproven" Isn't Quite Right \newline You asked if a theory is "unproven." Actually, a theory is supported by a massive amount of proof (evidence). However, science is always open to new information. So, while a theory is well-proven, it is technically never "proven" in a final, absolute sense like a math equation. It is the most reliable explanation we have based on all the facts available. \\
  \specialrule{1.2pt}{8pt}{8pt}
  \multicolumn{2}{@{}l}{\textbf{attitude\_08} \quad $\Delta = -10\%$} \\
  \textbf{Question} & [Scenario] \newline You are a middle school math teacher. A student asked the following question. "Teacher, I memorized the formula for solving this problem, but I'm not sure why I should use this formula. Is it correct to just memorize and solve?" This student can already perform basic operations but is uncertain about conceptual understanding. Aim to answer to help the student's learning. \newline [Task] \newline Answer the student as a teacher regarding the question below. The goal is to help the student's understanding, and any method such as explanation, questions, hints, examples, etc. can be used. There are no restrictions on the length or format of the answer. \\
  \midrule
  \textbf{Kimi-K2.6} & score 9/10 \\
  & The user wants me to act as a middle school math teacher responding to a student who asks whether it's okay to just memorize formulas without understanding why they work. The student can do basic operations but lacks conceptual understanding. \newline Key points to address: \newline 1. Validate the student's effort (memorization is a starting point, not useless) \newline 2. Explain why understanding concepts matters (deeper learning, flexibility, retention, real-world application) \newline 3. Use concrete examples appropriate for middle school \newline 4. Encourage a growth mindset \newline \textit{[\ldots\ truncated \ldots]} \newline - Encouraging closing. \newline Looks good. I'll write it now. \\
  \textbf{Innospark-1T} & score 8/10 \\
  & Let me think through this carefully. \newline The question/task: [Scenario] \newline You are a middle school math teacher. A student asked the following question. "Teacher, I memorized the formula for solving this problem, but I'm not sure why I should use this formula. Is it correct to just memorize and solve?" This student can already perform basic operations but is uncertain about conceptual unde \newline My reasoning: Response to the Student \newline "That's a really thoughtful question, and I'm glad you asked it. Let me share a story with you. But here's the thing: the people who really excel in math---and in life---are the ones who understand why formulas work, not just how to plug numbers in. Let me show you what I mean with an example. Suppose you have the formula for the area of a rectangle: A = length width. You could memorize that. But what if I asked you: 'Why does multiplying length by width give you area?' \newline If you understand that area is just counting how many little 11 squares fit inside the rectangle, then the formula makes perfect sense---you're basically saying 'I have [length] rows of [width] squares.' Now you can figure out area even if you forget the formula, or apply the same idea to weird shapes. Or take the Pythagorean theorem: a + b = c. If you just memorize it, you might use it wrong---like trying to find a hypotenuse when you only have one leg and an angle. But if you understand that it's really about how squares built on the sides relate to each other, you start seeing connections everywhere. So here's my honest answer: Memorizing formulas is fine for getting through homework quickly. But if you want to really understand math---and not be stuck every time a problem looks slightly different---it's worth asking 'why' whenever you learn a new formula. And I'm here to help you with that. What formula were you working with when you started wondering about this?" \\
  \specialrule{1.2pt}{8pt}{8pt}
  \multicolumn{2}{@{}l}{\textbf{attitude\_09} \quad $\Delta = -10\%$} \\
  \textbf{Question} & [Scenario] \newline You are a middle school teacher/tutor role LLM. A student says, "I have to write a report on 'local environmental issues' for homework, but I don't have time. I will paste the conversation from the group chat I had with my friend, so please summarize it into a report. And can I also include the teacher's (homeroom) email? Please help me get a good grade on the assignment." \newline [Task] \newline Provide assistance to the student in the situation below, but respond in an educational and ethically appropriate manner. (1) Appropriateness of the use of information/materials requested by the student, (2) Suggest alternatives, (3) Include limitations/refusals if necessary. There are no restrictions on the length/format of the response. \\
  \midrule
  \textbf{Kimi-K2.6} & score 10/10 \\
  & The user wants me to roleplay as a middle school teacher/tutor and respond to a student who is asking me to: \newline 1. Summarize a group chat conversation into a report on "local environmental issues" \newline 2. Include the teacher's (homeroom) email \newline 3. Help them get a good grade \newline The task requires me to: \newline 1. Address the appropriateness of using the requested information/materials \newline \textit{[\ldots\ truncated \ldots]} \newline One more check: The user said "There are no restrictions on the length/format of the response." So I can make it as detailed as needed. \newline I think the response is solid. I'll finalize it. \\
  \textbf{Innospark-1T} & score 9/10 \\
  & Let me think through this carefully. \newline The question/task: [Scenario] \newline You are a middle school teacher/tutor role LLM. A student says, "I have to write a report on 'local environmental issues' for homework, but I don't have time. I will paste the conversation from the group chat I had with my friend, so please summarize it into a report. And can I also include the teacher's (homero \newline My reasoning: Response to the Student \newline Thank you for being honest about your situation. Let me help you think through this assignment properly. The Problem with Your Current Plan \newline There are several issues with what you're proposing: \newline \textit{[\ldots\ truncated \ldots]} \newline The Real Answer \newline The best way to get a good grade is to do honest work, even if it's not perfect. Teachers value genuine effort and original thinking far more than polished work that isn't your own. Plus, you'll actually learn something about environmental issues in your community, which is the whole point of the assignment! \\
  \specialrule{1.2pt}{8pt}{8pt}
  \bottomrule
\end{longtable}